\scanpage{23}
\uline{Motifs} % L1

1) À tout préschéma \textsuperscript{noeth} (éventuellement de type fini sur un anneau % L2
noethérien) $X$ est associé une \uline{catégorie abélienne} % L3
$\mathcal M^+(X)$, dite catégorie des \uline{motifs} \inserted{\uline{effectifs}} sur $X$. C'est une % L4
$\mathbb Q$-catégorie \deleted{de $\mathbb Q$-vectoriels} abélienne, i.e. pour tt % L5
$M\in\mathrm{Ob}\,\mathcal M^+(X)$, et $n\in\mathbb Z$, $n\ne0$, $n\,\mathrm{id}_M$ est un isom % L6
de $M$. De plus, $\mathcal M^+(X)$ est muni d'un % L7
produit tensoriel commutatif et unitaire, \inserted{exact à droite,} % L8
% slanted
\margin{On peut en termes du \unsure{dernier|premier|données} construire des foncteurs $\unsure{\lambda^i|\bigwedge^i}M$ etc...} % L8a
l'unité est notée $\mathbb 1_X$ ou $\mathbb Q_X(0)$. % L9

On considère aussi la catégorie dérivée bornée % L10
$D^b(\mathcal M^+(X))$ de $\mathcal M^+(X)$. Le produit tensoriel est % L11
étendu en un bifoncteur $M\mathbin{\underline{\underline{\otimes}}}N$ en $M,N\in\mathrm{Ob}\,D^b(\mathcal M^+(X))$ % L12

2) Soit $f:X\longrightarrow Y$ un morphisme de préschémas % L13
\deleted{\unsure{loc}} noeth, il lui est associé un \uline{foncteur} % L14
\uline{exact} % L15
\[
f^*:\mathcal M^+(Y)\longrightarrow\mathcal M^+(X) % L16
\]
compatible avec $\otimes$, \inserted{d'où $\mathbb L f^*:D^{b}(\mathcal M^+(Y))\unsure{-|\to}D^{b}(\mathcal M^+(X))$} On a transitivité. % L17

Si $f$ est de type fini, et propre \uline{ou} $Y$ excellent, % L18
\unsure{On|on|on a} \deleted{\unsure{définit|défini}} de $=$ un foncteur % L19
\[
\mathbb R f_*:D^b(\mathcal M^+(X))\unsure{-|\longrightarrow}D^b(\mathcal M^+(Y)) % L20
\]

% archive sheet number: 22
